\documentclass[UTF8,zihao=-4,openany]{ctexbook}

\usepackage[a4paper,margin=2.45cm,headheight=15pt]{geometry}
\usepackage{amsmath,amssymb,bm,mathtools}
\usepackage{graphicx}
\usepackage{booktabs,tabularx,array}
\graphicspath{{figures/}}
\newcommand{\bookfig}[4]{%
  \begin{center}
    \includegraphics[width=0.92\linewidth,height=0.42\textheight,keepaspectratio]{#1}%
  \end{center}
  \noindent{\heiti 图~#2}\quad #3\par
  \vspace{0.35em}
  \noindent 符号：#4\par
  \vspace{0.6em}%
}
\newcommand{\bookfigbig}[4]{%
  \begin{center}
    \includegraphics[width=0.96\linewidth,height=0.58\textheight,keepaspectratio]{#1}%
  \end{center}
  \noindent{\heiti 图~#2}\quad #3\par
  \vspace{0.35em}
  \noindent 符号：#4\par
  \vspace{0.6em}%
}
\usepackage[dvipsnames]{xcolor}
\usepackage{fancyhdr}
\usepackage{enumitem}
\usepackage{listings}
\usepackage[
  unicode,
  colorlinks=true,
  linkcolor=Primary,
  urlcolor=Accent,
  bookmarksopen=true,
  bookmarksnumbered=true,
  bookmarksdepth=1
]{hyperref}

\definecolor{Primary}{HTML}{355C7D}
\definecolor{Accent}{HTML}{6C8EAD}
\definecolor{CodeBg}{HTML}{F6F7F8}

\setlength{\parindent}{2em}
\setlength{\parskip}{0.35em}
\setlist[itemize]{leftmargin=2em,itemsep=0.2em,topsep=0.3em}
\setlist[enumerate]{leftmargin=2.2em,itemsep=0.2em,topsep=0.3em}
\renewcommand{\arraystretch}{1.25}
\setcounter{secnumdepth}{1}
\setcounter{tocdepth}{1}
\allowdisplaybreaks[4]
\emergencystretch=3em
\hfuzz=5pt
\vfuzz=5pt

\pagestyle{fancy}
\fancyhf{}
\fancyhead[L]{计算方法与优化}
\fancyhead[R]{\hyperlink{contents}{返回目录}}
\fancyfoot[C]{\thepage}

\hypersetup{
  pdftitle={计算方法与优化学习笔记},
  pdfauthor={learning-notes}
}

\lstset{
  basicstyle=\ttfamily\small,
  backgroundcolor=\color{CodeBg},
  frame=single,
  rulecolor=\color{Accent},
  breaklines=true,
  columns=fullflexible,
  keepspaces=true,
  showspaces=false,
  showstringspaces=false,
  language=Python
}

\newcommand{\code}[1]{\texttt{#1}}
\renewcommand{\thesection}{第\arabic{section}页}

\title{\heiti 计算方法与优化\\[0.4em]\Large 学习笔记}
\author{内容提炼}
\date{\today}

\begin{document}

\maketitle
\frontmatter

\chapter*{说明}
\addcontentsline{toc}{chapter}{说明}
本笔记对应武汉大学数学与统计学院戴书洋老师
\emph{计算方法与优化}（Computational Methods and Optimization）
2026--2027~秋季课程第~1~讲课件
\texttt{CM+Opt\_ch1.pdf}（共~57~页）。
每一页课件对应正文一个小节，节号写成「第$n$页」，标题沿用该页课件标题；
同名页在标题后用括号标出该页特有主题，不改课件原词。
试作版只充实当前这一份课件；后续各讲待提供后再按同样规则追加，此处不编造目录。
正文以公式陈述定义、成立条件、关键推导与结论。
每一公式后写明：数学上如何由输入算出输出，以及得到的数在问题中的实际含义。
有信息量的示意图按课件页码编号插入；表按课件行列重排。

\hypertarget{contents}{}
\tableofcontents

\mainmatter
\chapter{课程导论}

本章对应课件第~1--20~页：课程信息、为什么需要计算数学、算法观念，以及学期内容安排。

\section{计算方法与优化}

课程名\emph{计算方法与优化}，学期为 2026--2027~秋季，主讲戴书洋，单位为武汉大学数学与统计学院。
这一页只给出身份信息，没有公式。
它规定后续每一页都在同一门课的记号系统下展开：先把连续模型写成有限步可执行的算法，再分析误差、稳定性和运算量。

\section{课程信息}

联系与成绩按课件给出：
邮箱 \code{shuyang\_dai@whu.edu.cn}，答疑用邮件或 QQ，QQ~群~337413555，办公室雷军科技楼~719。
总评
\[
\text{总评}=0.6\times\text{期末}+0.4\times\text{平时}.
\]
输入是百分制的期末分与平时分，输出是加权标量，取值通常在 $[0,100]$。
$0.6$ 与 $0.4$ 不是物理量，只是成绩权重；期末权重大，表示课程仍以可独立完成的书面与计算考查为主。

\section{大模型这么强，为什么还要上（数学、编程）课？}

课件用大模型对话截图提出问题：生成式模型已经能写代码、解算例，为何仍要学数学与编程。
这一页没有给出公式，答案在后续页展开：模型把问题翻译成方程之后，真正决定对错的是算法是否收敛、误差如何传播、浮点运算能否稳住。
大模型输出的是文本；本课要的是可验证的数值，取值范围由误差限与稳定性定理约束，不能用「看起来合理」代替。

\section{纯粹数学与应用数学}

\bookfig{fig-p04-math-split.png}{第4页}{数学分成纯粹数学与应用数学两条路线。}{纯粹数学：问题来自数学内部；应用数学：问题来自科学与工程，强调建模加求解}

课件把数学分成两支。
纯粹数学不以应用为目的，强调严格、抽象与结构，问题如哥德巴赫猜想、黎曼猜想；代表人举华罗庚、陈省身、丘成桐、陈景润、张益唐。
应用数学关心科学应用，路线是
\[
\text{实际对象}\to\text{数学模型}\to\text{数值问题}\to\text{算法与程序}.
\]
输入是物理、工程或数据中的对象，输出是可计算的方程与算法；中间每一步都丢掉一部分细节，因此后面必须估计误差。
问题来源举流体力学、材料科学、神经科学；代表人举冯·诺依曼、香农、维纳、纳什、华罗庚。
本课落在应用数学一侧：同一位华罗庚同时出现在两支，说明「算得出来」与「证得出来」并不互斥。

\section{数学重要性：以计算为例（计算古已有之）}

计算并非计算机发明后才出现。
天文用计算研究星球运动与历法，军事用计算分析弹道，航海用计算选择航线。
课件用第谷收集观测、开普勒从数据中找规律作为代表：
观测给出有限个点 $\{(t_i,x_i)\}$，规律是一个待定函数 $x(t)$。
输入是离散数据，输出是可外推的模型；数据有观测误差，所以后文必须讨论拟合与插值，而不是假设点落在精确曲线上。

\section{数学重要性：以计算为例（现代计算自二战开始）}

课件把现代大规模计算的起点放在第二次世界大战：研制核武器需要流体与辐射的数值模拟，破译密码需要对离散结构做穷举与统计。
这一页的图是历史场景，不进入公式。
它要说明的数量级变化是：问题规模一旦超出手算，就必须把连续方程离散化，并接受截断误差与舍入误差。

\section{数学重要性：以计算为例（计算必须算到目标精度）}

课件用课堂板书场景强调：理论公式写完之后，仍要把数值「算到该停的地方」。
对一个级数或迭代，停止准则是误差降到事先给定的容许值 $\varepsilon>0$，而不是「公式已经写出来」。
输入是算法与 $\varepsilon$，输出是满足
\[
\lvert\text{近似}-\text{真值}\rvert\le\varepsilon
\]
的有限步结果；$\varepsilon$ 的量纲与未知量相同，取值由应用精度决定。

\section{数学重要性：以计算为例（自己动手从零开始）}

课件用大规模手算的历史画面说明：没有现成黑箱时，必须从最基本的四则运算搭出算法。
后文秦九韶算法、浮点舍入、范数与条件数，都是「从零开始」的具体化：先规定运算顺序，再计算运算次数，再估计误差放大倍数。

\section{计算的变革}

电子计算机把计算从手算变成自动的有限步算术。
ENIAC 于 1946~年~2~月~14~日在美国投入使用，课件称其为世界上第一台通用计算机，占地约 $170$~平方米，重约 $30$~英吨，速度为每秒 $5000$~次加法或 $400$~次乘法。
这些数的输入是历史记载，输出是量级对照：加法比乘法大约快一个数量级，因此后文统计运算量时要把乘法次数单独列出。
计算机只改变执行速度，不取消误差：每一步仍是有限位算术。

\section{计算的用途：核武器数值模拟}

课件用核装置数值模拟作为计算的用途之一。
控制方程是连续的偏微分方程，机器只能解有限维代数问题，因此必须离散。
输入是连续场，输出是网格上的有限个未知数；网格越细，截断误差通常越小，未知数个数与运算量同时上升。
这一页的照片只说明应用场景，不提供可抄的公式。

\section{计算的用途：大飞机制造}

大飞机的气动、结构与控制同样依赖数值模拟。
与上一页相同：把偏微分方程变成大规模线性或非线性方程组。
输入是几何与工况，输出是可制造、可试验对照的数值预测；预测值的可信区间由网格、模型与舍入共同决定。

\section{计算能力的驱动力}

课件的结论是：计算能力的驱动来自需求，而不是先有机器再找问题。
国家需要核武器、大飞机、天气预报等能力时，受政治、经济与时效约束，只能靠计算在试验之前给出可检验的数量。
需求规定容许误差与时限，时限再反过来规定算法的运算量上限。

\section{基本问题}

计算方法要回答的基本问题是：
如何把数学模型写成数值问题；如何确定算法的适用范围；如何估计给定算法的精度；误差如何在计算中积累和传播；如何构造精度更高的算法；如何减少存储；如何分析运算量与运算时间。
课件把本课内容分成三块：
数值线性代数（线性方程组、特征值与特征向量），
数值逼近（插值与逼近），
优化理论与算法（如何训练神经网络）。
这三块对应后文的 $Ax=b$、用简单函数代替复杂函数、以及
\[
\min_{\theta}\,L(\theta).
\]
最后式的输入是参数 $\theta$ 与训练损失 $L$，输出是标量损失值；训练要的是使 $L$ 下降的 $\theta$，而不是把 $L$ 的表达式写得更长。

\section{学习方法}

核心是算法的原理及实现。
基本任务：理解每个算法的数学背景、原理和线索，对最基本的算法非常熟悉。
理论部分：认识算法，并能做收敛性、稳定性分析。
实现部分：各章算法的目的是实际计算，必须能上机并分析结果。
因此后文每个公式都要能回答两件事：用哪些输入算出什么输出，以及这个数在原问题里代表误差、残差、步长还是条件数。

\section{算法}

算法不是单独一条数学公式，而是由基本运算及其顺序组成的完整求解方案。
可用流程图描述；选择算法是数值计算的关键环节，不同算法有不同的优势与劣势。
若记基本运算集合为 $\mathcal{O}$、顺序为 $\prec$，则一个算法是 $(\mathcal{O},\prec)$。
输入是问题实例，输出是有限步之后的近似解；运算次数是非负整数，运行时间与它成正比，但不等于墙钟时间。

\section{算法（多项式与秦九韶）}

多项式
\[
P_n(x)=a_n x^n+a_{n-1}x^{n-1}+\cdots+a_1 x+a_0
\]
的输入是系数 $(a_0,\ldots,a_n)$ 与求值点 $x$，输出是标量 $P_n(x)$。
$P_n(x)$ 就是该多项式在 $x$ 处的函数值，量纲与 $a_0$ 相同。

方法一：先算每个 $a_i x^i$（$i=0,1,\ldots,n$）再相加。
乘法次数为
\[
1+2+\cdots+n=\frac{n(n+1)}{2},
\]
加法次数为 $n$。
左边对 $i=1,\ldots,n$ 统计「$x$ 连乘 $i$ 次」的乘法；输出是次数，单位是次，取值 $\{0,1,2,\ldots\}$。

方法二：秦九韶算法（1247）/Horner 算法（1819），
\[
P_n(x)=\bigl(\cdots((a_n x+a_{n-1})x+a_{n-2})x+\cdots+a_1\bigr)x+a_0.
\]
递归为 $b_n=a_n$，
\[
b_i=b_{i+1}x+a_i,\qquad i=n-1,n-2,\ldots,0,
\]
且 $P_n(x)=b_0$。
每步输入已有 $b_{i+1}$、同一个 $x$ 和系数 $a_i$，输出新的标量 $b_i$；
$b_i$ 是从最高次项看到第 $i$ 次的嵌套值，最后的 $b_0$ 才是多项式值。
乘法次数为 $n$，加法次数为 $n$，比方法一少一个数量级的乘法。

\begin{lstlisting}
def horner(a, x):
    # a[0] = a_0, ..., a[n] = a_n
    b = a[-1]
    for ai in reversed(a[:-1]):
        b = b * x + ai
    return b
\end{lstlisting}
对 $P(x)=2x^2+3x+4$ 即 \code{a=[4,3,2]}，\code{horner(a,5)} 得到 $69$，与 $2\cdot 25+15+4$ 相同；返回值就是 $P(5)$，不是系数。

\section{算法（Basel 问题的直接求和）}

Basel 问题为
\[
\sum_{k=1}^{\infty}\frac{1}{k^2}=\frac{\pi^2}{6}.
\]
用部分和
\[
P_n=\sum_{k=1}^{n}\frac{1}{k^2}
\]
逼近 $\pi^2/6$。
输入是截断下标 $n\in\mathbb{N}$，输出是正标量 $P_n\in(0,\pi^2/6)$；
$P_n$ 是前 $n$ 项的数值和，不是 $\pi$ 本身。

直接计算的截断误差为
\[
\frac{\pi^2}{6}-P_n=\sum_{k=n+1}^{\infty}\frac{1}{k^2}.
\]
由
\[
\frac{1}{k(k+1)}<\frac{1}{k^2}<\frac{1}{k(k-1)}\quad(k\ge n+1)
\]
得到
\[
\frac{1}{n+1}<\sum_{k=n+1}^{\infty}\frac{1}{k^2}<\frac{1}{n}.
\]
左右两端是误差的可计算上下界，量纲与 $1/k^2$ 相同，取值随 $n$ 按 $O(n^{-1})$ 下降。
因此要把误差压到 $10^{-2m}$，大致需要 $n\sim 10^{2m}$ 项，直接求和很慢。

\bookfig{fig-p17-basel-direct.png}{第17页}{直接部分和的误差（左）与对数误差（右）。}{$n$：项数；Error：$\pi^2/6-P_n$；Log Error：误差的常用对数；虚线：课件给出的 $1/n$、$1/(n+1)$ 界}

\section{算法（Basel 问题的加速求和）}

利用
\[
\sum_{k=1}^{\infty}\frac{1}{k(k+1)}=1,
\qquad
\frac{1}{k^2}-\frac{1}{k(k+1)}=\frac{1}{k^2(k+1)},
\]
改写
\[
\sum_{k=1}^{\infty}\frac{1}{k^2}
=1+\sum_{k=1}^{\infty}\frac{1}{k^2(k+1)}
\approx 1+\sum_{k=1}^{n}\frac{1}{k^2(k+1)}.
\]
新部分和的输入仍是 $n$，输出是
\[
Q_n=1+\sum_{k=1}^{n}\frac{1}{k^2(k+1)},
\]
它仍是 $\pi^2/6$ 的近似，但余项满足
\[
\sum_{k=n+1}^{\infty}\frac{1}{k^2(k+1)}
<\frac{1}{n}\sum_{k=n+1}^{\infty}\frac{1}{k(k+1)}
<\frac{1}{n^2}.
\]
$1/n^2$ 是余项上界，比直接法的 $1/n$ 高一阶。
同一 $n$ 下 $Q_n$ 比 $P_n$ 更接近 $\pi^2/6$；对数误差图上加速后的点下降明显更快。

\bookfig{fig-p18-basel-accel.png}{第18页}{加速后的误差（左）与对数误差（右），绿色点为加速方案。}{$n$：项数；蓝点：直接法；绿点：加 $1/(k^2(k+1))$ 后的余项}

\section{算法（算法优劣的影响）}

算法优劣直接影响速度和效率。
对小型问题，运算次数与内存似乎无所谓；对大型问题，它们起决定作用。
不合适的算法还会因为近似性和误差的传播、积累而损害精度，甚至使计算失败。
结合前两页：同样逼近 $\pi^2/6$，把余项从 $O(n^{-1})$ 改到 $O(n^{-2})$ 并不改问题，只改算法；
输出的误差数量级可以差几个 $10$ 的幂。

\section{上课内容}

学期按周安排为：
第~1~周基础；
第~2--4~周线性方程的直接法、迭代法、梯度法（最速下降与共轭梯度）；
第~5--6~周最小二乘；
第~7--8~周非线性方程与 SVD；
第~9--10~周插值逼近；
第~11--15~周优化。
这一页是后续各讲的课表，不是第~1~章内部小节名。
当前笔记只把「第~1~周基础」写完；其余周次等对应课件后再按页追加。

\chapter{基础知识}

课件把本章分成五块：介绍，误差，浮点数与舍入误差，范数，线性方程组的性态与算法稳定性。
下面按同一顺序，每块先写对象与公式，再给一两个算例说明公式里的数代表什么。

\section{介绍}

用计算机解题的固定链条是
\[
\text{模型}\to\text{数值问题}\to\text{算法}\to\text{程序}.
\]
输出始终是近似值，因此后文只关心三件事：收敛、稳定、运算量。
本课遇到的数值问题可以写成三类：
\begin{align}
Ax&=b,\\
f(x)&=0,\\
\min_{\theta}\,L(\theta).
\end{align}
第一式输入 $A\in\mathbb{R}^{n\times n}$、$b\in\mathbb{R}^n$，输出 $x\in\mathbb{R}^n$，是未知量本身。
第二式输入非线性 $f$ 与初值，输出根 $x$。
第三式输入参数 $\theta$ 与损失 $L$，输出使 $L$ 变小的 $\theta$；$L(\theta)$ 是非负标量风险，不是预测值。

成立条件随问题而变：方形可逆组有唯一解；超定组 $m>n$ 一般无精确解，改成
\[
\min_x\,\lVert Ax-b\rVert_2^2
\quad\Leftrightarrow\quad
A^{\mathsf T}Ax=A^{\mathsf T}b;
\]
非线性与训练还要算法收敛。
形式上写 $x=A^{-1}b$ 并不要求真去求逆：求逆与乘法同为 $O(n^3)$，常数更大，实际用 LU、Cholesky、QR 或迭代直接得到一个 $x$。

\textbf{案例 A（线性组）。}
《九章算术》
\[
\begin{pmatrix}3&2&1\\2&3&1\\1&2&3\end{pmatrix}
\begin{pmatrix}x\\y\\z\end{pmatrix}
=\begin{pmatrix}39\\34\\26\end{pmatrix}
\]
的解是 $x=9.25$，$y=4.25$，$z=2.75$，单位为斗/秉。
这是 $Ax=b$ 的原型。

\textbf{案例 B（非线性）。}
Kepler 方程 $x-\varepsilon\sin x-t=0$（$0<\varepsilon<1$）对 $x$ 非线性，$x(t)$ 是偏近点角。
水温插值与人口拟合则分别是「穿过已知点」和「点多于参数时的最小二乘」：
人口用 $y\approx\beta_1 t^3+\beta_2 t^2+\beta_3 t+\beta_4$，输入 $\{(t_i,y_i)\}$，输出系数 $\beta$，再代入新 $t$ 得到预测人口。
神经网络把 $f_\theta$ 取得更宽，训练仍是 $\min_\theta L(\theta)$，梯度步
\[
\theta_{k+1}=\theta_k-\eta\nabla L(\theta_k)
\]
的 $\eta>0$ 是步长，$\nabla L$ 是损失对参数的梯度。

\textbf{案例 C（运算量）。}
朴素矩阵乘法 $(AB)_{ij}=\sum_{k=1}^n a_{ik}b_{kj}$ 要 $n^3$ 次乘法；
$n=10^5$、每秒 $10^9$ 次时约 $11.6$ 天。
Strassen 把指数从 $3$ 降到 $\log_2 7\approx 2.807$，说明同一问题可以换算法，但工程上仍受稳定性与缓存制约。

\section{误差}

数值分析假定模型与观测已给定，只跟踪两类计算误差：截断（方法）误差与舍入误差。
设 $x^*$ 为真值，$x$ 为近似值。绝对误差与相对误差为
\begin{equation}
\varepsilon=x-x^*,
\qquad
\xi=\frac{x-x^*}{x^*}.
\end{equation}
$\varepsilon$ 与 $x^*$ 同量纲，可正可负；$\lvert\varepsilon\rvert\le\varepsilon^*$ 时称 $\varepsilon^*$ 为绝对误差限，常写 $x^*=x\pm\varepsilon^*$。
$\xi$ 无量纲，$\lvert\xi\rvert=10^{-k}$ 大致对应 $k$ 位有效数字。
真值未知时改用 $\lvert x^*-x\rvert/\lvert x\rvert$。
$\lvert\varepsilon\rvert$ 小并不等于「够准」：同样差 $1$，对 $10^6$ 和对 $1$ 的意义不同，所以比较精度看 $\xi$。

有效数字：写 $x=\pm 0.a_1\ldots a_n\times 10^m$（$a_1\neq 0$），若
\[
\lvert x-x^*\rvert\le 0.5\times 10^{m-n},
\]
则 $x$ 有 $n$ 位有效数字。$n$ 是位数，不是 $x$ 的值。

误差沿函数传播。一元
\[
\Delta f(x^*)\approx\lvert f'(x^*)\rvert\,\lvert x-x^*\rvert,
\]
$\lvert f'(x^*)\rvert$ 是绝对误差的局部放大率。
二元一阶上界为
\[
\Delta f\approx\bigl\lvert\partial f/\partial u\bigr\rvert\Delta u
+\bigl\lvert\partial f/\partial v\bigr\rvert\Delta v.
\]
四则运算的估计误差：加减都是 $\Delta\tilde u+\Delta\tilde v$；
乘为 $\lvert\tilde u\rvert\Delta\tilde v+\lvert\tilde v\rvert\Delta\tilde u$；
除为同一分子除以 $\lvert\tilde v\rvert^2$。
减法的绝对误差限与加法相同，但真值接近时相对误差会被放大（相消）。

\bookfig{fig-p41-error-propagation.png}{第41页}{切线把自变量误差变成函数值误差的一阶估计。}{$\Delta\tilde x$：自变量绝对误差；估计误差：$\lvert f'\rvert\Delta\tilde x$；真误差：$\lvert f(x)-f(\tilde x)\rvert$}

\textbf{案例 A（有效数字）。}
$x^*=\pi$，$x=3.1415=0.31415\times 10^1$，故 $m=1$。
$\lvert\pi-3.1415\rvert\approx 9.3\times 10^{-5}$ 大于 $0.5\times 10^{1-5}$、小于 $0.5\times 10^{1-4}$，因此只有 $4$ 位有效数字。
四舍五入到四位：$42.195\to 42.20$，$0.0375551\to 0.03756$。

\textbf{案例 B（截断）。}
$I=\int_0^1 e^{-x^2}\,\mathrm{d}x$ 无初等原函数。逐项积分
\[
I=\sum_{k=0}^{\infty}\frac{(-1)^k}{k!\,(2k+1)},
\qquad
I_n=\sum_{k=0}^{n}\frac{(-1)^k}{k!\,(2k+1)},
\]
交错余项
\[
\lvert I-I_n\rvert\le\frac{1}{(n+1)!\,(2n+3)}.
\]
$I_n$ 是 $I$ 的近似（真值 $\approx 0.746824$）；右端是可事先计算的截断上界，不是随机误差。
$n=4$ 时 $I_4\approx 0.74749$，上界 $\approx 7.6\times 10^{-4}$，真误差 $\approx 6.6\times 10^{-4}$。

\section{浮点数及舍入误差}

浮点系统 $F(\beta,t,m,M)$ 表示
\[
x=W\beta^j,
\qquad
W=\pm\sum_{r=1}^{t}d_r\beta^{-r},
\quad -m\le j\le M.
\]
$\beta$ 是基，$t$ 是尾数位，$j$ 是阶码。规格化要求 $d_1\neq 0$（二进制即 $d_1=1$）。
$F$ 是有限散点集，不是 $\mathbb{R}$。
IEEE 单精度 $1+8+23$ 位，双精度 $1+11+52$ 位，数值为 $(-1)^S W 2^j$。

相邻规格化数间距 $\Delta=\beta^{j-t}$。四舍五入到最近浮点 $x_f$ 时
\[
\lvert x_f-x_R\rvert\le\tfrac12\beta^{j-t}.
\]
结合规格化下界 $\lvert x_R\rvert\ge\beta^{j-1}$，得到定理~1.3.1：
\begin{equation}
x_f=x_R(1+\delta),
\qquad
\lvert\delta\rvert\le u:=\tfrac12\beta^{1-t}.
\end{equation}
$\delta$ 是相对舍入，无量纲；$u$ 是单位舍入，只依赖 $\beta,t$，与 $\lvert x_R\rvert$ 无关（不溢出、规格化时）。
二进制 $u=2^{-t}$；双精度 $t=53$（含隐藏位）时 $u=2^{-53}\approx 1.11\times 10^{-16}$。

\bookfig{fig-p46-rounding-proof.png}{第46页}{由间距与规格化下界推出相对误差界 $u$。}{$\Delta$：相邻浮点间距；$u$：单位舍入；$\delta$：相对误差}

\section{向量范数与矩阵范数}

$\mathbb{R}^n$ 上的范数满足正定、绝对齐次、三角不等式。常用
\begin{align}
\lVert x\rVert_1&=\sum_i\lvert x_i\rvert,\\
\lVert x\rVert_2&=\bigl(\sum_i x_i^2\bigr)^{1/2},\\
\lVert x\rVert_\infty&=\max_i\lvert x_i\rvert.
\end{align}
输入向量，输出非负标量，表示长度。
有限维上一切范数等价：存在 $c_1,c_2>0$ 使 $c_1\lVert x\rVert'\le\lVert x\rVert\le c_2\lVert x\rVert'$，
因此「$x_k\to x^*$」与「$\lVert x_k-x^*\rVert\to 0$」对任意范数同时成立。
$\lVert x_k-x^*\rVert$ 是误差长度，不是某一个坐标。

矩阵范数另要求相容 $\lVert AB\rVert\le\lVert A\rVert\lVert B\rVert$。
Frobenius 范数把 $A$ 当成向量做 $2$-范数。
由向量范数诱导的算子范数
\[
\lVert A\rVert=\max_{x\neq 0}\frac{\lVert Ax\rVert}{\lVert x\rVert}
\]
是单位球上的最大拉伸，也是与该向量范数相容的最小矩阵范数，且 $\lVert I\rVert=1$。
列和、行和与谱范数分别为
\begin{align}
\lVert A\rVert_1&=\max_j\sum_i\lvert a_{ij}\rvert,\\
\lVert A\rVert_\infty&=\max_i\sum_j\lvert a_{ij}\rvert,\\
\lVert A\rVert_2&=\sqrt{\rho(A^{\mathsf T}A)}.
\end{align}
谱半径 $\rho(A)\le\lVert A\rVert$ 对任意矩阵范数成立；$\rho(A)$ 是按模最大特征值，控制 $A^k$ 是否趋于零。
若某算子范数下 $\lVert A\rVert<1$，则 $I\pm A$ 可逆，且
\[
\frac{1}{1+\lVert A\rVert}\le\lVert(I\pm A)^{-1}\rVert\le\frac{1}{1-\lVert A\rVert}.
\]

\textbf{案例。}
$x=(1,-2,3)^{\mathsf T}$：$\lVert x\rVert_1=6$，$\lVert x\rVert_2=\sqrt{14}$，$\lVert x\rVert_\infty=3$。
$A=\begin{pmatrix}1&-2\\-3&4\end{pmatrix}$：
$\lVert A\rVert_1=6$，$\lVert A\rVert_\infty=7$，
$\lVert A\rVert_2=\sqrt{15+\sqrt{221}}$。
对任意算子范数还有 $\lVert A^{-1}\rVert\ge 1/\lVert A\rVert$，以及
$\lVert A^{-1}-B^{-1}\rVert\le\lVert A^{-1}\rVert\lVert B^{-1}\rVert\lVert A-B\rVert$。

\section{线性方程组的性态，算法的稳定性}

解 $Ax=b$。右端扰动 $\delta b$ 给出
\begin{equation}
\frac{\lVert\delta x\rVert}{\lVert x\rVert}
\le
\operatorname{cond}(A)\,\frac{\lVert\delta b\rVert}{\lVert b\rVert},
\qquad
\operatorname{cond}(A)=\lVert A\rVert\lVert A^{-1}\rVert.
\end{equation}
$\operatorname{cond}(A)\ge 1$，无量纲，$\operatorname{cond}(kA)=\operatorname{cond}(A)$（$k\neq 0$）。
它是相对误差放大因子：输入相对右端扰动，输出相对解扰动的上界。
矩阵扰动 $\delta A$ 时还多一个分母 $1-\operatorname{cond}(A)\lVert\delta A\rVert/\lVert A\rVert$（需为正）。
$\operatorname{cond}(A)\gg 1$ 称病态，反之为良态。
经验上，行列式极小、行近似相关、主元过小、特征值数量级悬殊，都提示病态。
算法稳定性说的是舍入在计算过程中被放大多少；病态是问题固有的，稳定算法也不能凭空变出有效位。

\textbf{案例 A。}
\[
\begin{cases}x_1+x_2=2,\\ x_1+1.001 x_2=2,\end{cases}
\quad
\begin{cases}x_1+x_2=2,\\ x_1+1.001 x_2=2.001.\end{cases}
\]
解从 $(2,0)$ 变成 $(1,1)$。右端相对扰动约 $5\times 10^{-4}$，解的相对变化为 $O(1)$，放大约 $10^3$，与该 $2\times 2$ 矩阵的条件数同量级。

\textbf{案例 B（Hilbert）。}
\[
H_{ij}=\frac{1}{i+j-1},\quad i,j=1,\ldots,15,
\qquad
b=H\mathbf{1}.
\]
精确解 $x^*=\mathbf{1}$。双精度 $\operatorname{cond}_2(H)\approx 8.72\times 10^{17}$，
$u\cdot\operatorname{cond}_2(H)=O(10^2)$，有效数字可以全部丢失。
浮点解的残差 $\lVert Hx-b\rVert_2/\lVert b\rVert_2\approx 2\times 10^{-16}$，
相对误差 $\lVert x-x^*\rVert_2/\lVert x^*\rVert_2\approx 4.93$；
分量可偏离到 $-8.7$、$12.8$。
残差达到机器精度并不表示解对，这就是病态。
\begin{lstlisting}
import numpy as np
n, i = 15, np.arange(1, 16)
H = 1.0 / (i[:, None] + i[None, :] - 1.0)
x = np.linalg.solve(H, H @ np.ones(n))
\end{lstlisting}
\code{solve} 输出满足 $Hx\approx b$ 的 $x\in\mathbb{R}^{15}$，它不是 $x^*$。


\end{document}
